\chapter{Complexity — Call Return}
%%% generated by the blueprint pipeline from TCSlib.Complexity.SpaceComplexity.Machines.CallReturn
\section{Overview}

Declarations of the Lean module \texttt{TCSlib.Complexity.SpaceComplexity.Machines.CallReturn}.

\section{Declarations}

\begin{definition}[Combined machine configuration]
\lean{Complexity.LogProg.mkCfg}
\label{Complexity.LogProg.mkCfg}
\leanok
\uses{Complexity.LogProg.CSt, Turing.Cfg}
Let $x$ be a finite Boolean input, let $m$ be the number of registers, and let $k_D$ be
the number of simulated machine work tapes. Given a control state, an input-head
position, register-tape contents and head positions, simulated work-tape contents and
head positions, and an output list, define the combined configuration to have the given
control state, input-head position, and output, with the register tapes and simulated
work tapes concatenated in that order, and likewise their head positions.
\end{definition}

\begin{lemma}[State of the constructed configuration]
\lean{Complexity.LogProg.mkCfg_state}
\label{Complexity.LogProg.mkCfg_state}
\leanok
\uses{Complexity.LogProg.CSt, Complexity.LogProg.mkCfg}
\difficulty{1}
For any optional compiled-machine state, valid input-head position, register-tape
contents and head positions, decider-tape contents and head positions, and output word,
the state of the configuration assembled from these data is the specified optional
state.
\end{lemma}

\begin{lemma}[Effect of register-only head moves]
\lean{Complexity.LogProg.apply_moves}
\label{Complexity.LogProg.apply_moves}
\leanok
\uses{Complexity.LogProg.CSt, Complexity.LogProg.dIdle, Complexity.LogProg.mkCfg, Turing.Action, Turing.Action.apply, Turing.moveInputPos}
\difficulty{4}
Let a configuration with input $x$ have state $st$, input-head position $ip$, register
tapes $rt$ and head positions $rp$, decider tapes $dt$ and head positions $dp$, and
output $out$. Given an input-head movement $mv_I$ and a movement sign $f(r)$ for each
register $r$, an action that makes those moves, leaves all tape contents unchanged,
emits no output, and enters state $st'$ sends the configuration to one with state $st'$,
input-head position moved by the clamped move $mv_I$, register-head positions
$rp(r)+f(r)$, and the decider tapes and all other listed data unchanged.
\end{lemma}

\begin{lemma}[Moving a register head]
\lean{Complexity.LogProg.apply_regMove}
\label{Complexity.LogProg.apply_regMove}
\leanok
\uses{Complexity.LogProg.CSt, Complexity.LogProg.apply_moves, Complexity.LogProg.dIdle, Complexity.LogProg.mkCfg, Complexity.LogProg.regMove, Turing.Action, Turing.Action.apply, Turing.moveInputPos_zero}
\difficulty{4}
Let $m,d,k_D\in\mathbb N$, let $\Lambda$ and $SD$ be state types, and let $x$ be a
Boolean input. For any optional control state, input-head position, register tapes and
head positions, decider tapes and head positions, output word, register
$r\in\mathrm{Fin}(m)$, movement $mv\in\{-1,0,1\}$, and new control state
$st'\in\mathrm{CSt}(\Lambda,SD,m)$, the action that leaves the input head and tapes
unchanged, moves register $r$ by $mv$, leaves all decider heads fixed, and enters $st'$
transforms the configuration into the one with state $st'$ and register-head position at
$r$ increased by $mv$, with all other specified data unchanged.
\end{lemma}

\begin{lemma}[Effect of an input-only action]
\lean{Complexity.LogProg.apply_inputOnly}
\label{Complexity.LogProg.apply_inputOnly}
\leanok
\uses{Complexity.LogProg.CSt, Complexity.LogProg.mkCfg, Turing.Action, Turing.Action.apply, Turing.moveInputPos}
\difficulty{4}
Let a configuration with input $x$ have state $st$, input-head position $ip$, register
tapes $rt$ and positions $rp$, decider tapes $dt$ and positions $dp$, and output $out$.
An action that moves the input head by $mv_I$, leaves every work tape unchanged, emits
no output, and enters state $st'$ sends this configuration to the one with state $st'$,
input-head position moved by the clamped move $mv_I$, and all work tapes, head
positions, and output unchanged.
\end{lemma}

\begin{lemma}[Seam configuration decomposition]
\lean{Complexity.LogProg.seam_eq_mkCfg}
\label{Complexity.LogProg.seam_eq_mkCfg}
\leanok
\uses{Complexity.LogProg.mkCfg, Complexity.LogProg.seam, Turing.Cfg}
\difficulty{1}
Let $c$ be a configuration with $m$ Boolean work tapes and state set $\Lambda$, and let
$k_D$ be a natural number. The configuration obtained by mapping its state into the
program-state component and adjoining $k_D$ blank work tapes with heads at position $0$
is equal to the configuration assembled from its original state, input-head position,
work tapes, work-head positions, and output together with those same $k_D$ blank tapes
and zero-position heads.
\end{lemma}

\begin{lemma}[Input rewind after a return]
\lean{Complexity.LogProg.ret2_run}
\label{Complexity.LogProg.ret2_run}
\leanok
\uses{Complexity.LogProg.CSt, Complexity.LogProg.CallSpec, Complexity.LogProg.apply_inputOnly, Complexity.LogProg.compileTM, Complexity.LogProg.ctr, Complexity.LogProg.mkCfg, Complexity.LogProg.mkCfg_state, Complexity.LogProg.nextRet, Turing.Cfg, Turing.Cfg.inputSymbol, Turing.FinTM.moveInputPos_neg_val, Turing.MultiTapeTM.runFrom, Turing.MultiTapeTM.runFrom_succ_eq_step, Turing.MultiTapeTM.step, Turing.inputSymbolInner, Turing.moveInputPos}
\difficulty{6}
Let $P$ be a register program with $m$ registers and $d$ deciders, let $D$ be a Boolean
multi-tape Turing machine with $k_D$ work tapes, and let $x$ be a Boolean input. Suppose
the call specification at node $l$ is $cs$. Starting from the compiled machine in its
second return phase at $l$, with segment index $s$, direction bit $dir$, result bit $b$,
arbitrary register and decider tapes, head positions, and output, take an input-head
position $ip$ whose value is $j$, where $j\le |x|$. Through every time $t\le j+1$, all
work-tape head positions remain unchanged; after $j+1$ steps the configuration has
input-head position $1$, the same tapes, work-head positions and output, and state given
by the return successor determined by $cs,l,s,dir,b$ and register index $0$.
\end{lemma}

\begin{lemma}[Completing a register scan]
\lean{Complexity.LogProg.retR_scan}
\label{Complexity.LogProg.retR_scan}
\leanok
\uses{Complexity.LogProg.CSt, Complexity.LogProg.CallSpec, Complexity.LogProg.apply_regMove, Complexity.LogProg.compileTM, Complexity.LogProg.ctr, Complexity.LogProg.mkCfg, Complexity.LogProg.mkCfg_state, Complexity.LogProg.nextRet, Turing.Cfg, Turing.Cfg.workTapeSymbols, Turing.FinTM.bufferTape, Turing.MultiTapeTM.runFrom, Turing.MultiTapeTM.runFrom_succ_eq_step, Turing.MultiTapeTM.step}
\difficulty{6}
Let $P$ be a register program with $m$ registers and $d$ deciders, let $D$ be a Boolean
multi-tape Turing machine with $k_D$ work tapes, and let $x$ be a Boolean input. Suppose
node $l$ has call specification $cs$, the register selected by argument position $a$ is
$rr$, and register $rr$ contains the buffer tape for a Boolean word $w$. If the machine
starts in the scanning return phase for this register and its head positions $rp$
satisfy $rp(rr)+1=n$ for a natural number $n\le |w|$, then at time $n+1$ it has advanced
to the return successor for argument position $a+1$, with register $rr$'s head at zero
and all other register head positions unchanged. At every time $t\le n+1$, each register
head either remains at its initial position or, for register $rr$, lies between $-1$ and
the maximum of its initial position and $0$; the decider tape head positions remain
unchanged throughout.
\end{lemma}

\begin{definition}[Register position on the left side]
\lean{Complexity.LogProg.leftSide}
\label{Complexity.LogProg.leftSide}
\leanok
Let $m$ be a natural number, let $s$ be an index in $\{0,\ldots,m\}$, let $a$ be a
register index in $\{0,\ldots,m-1\}$, and let $\mathrm{dir}$ be a Boolean. Define the
left-side indicator to be true exactly when $s<a+1$, or when $s=a+1$ and $\mathrm{dir}$
is false.
\end{definition}

\begin{lemma}[Restoring a register head]
\lean{Complexity.LogProg.retR_run}
\label{Complexity.LogProg.retR_run}
\leanok
\uses{Complexity.LogProg.CSt, Complexity.LogProg.CallSpec, Complexity.LogProg.apply_regMove, Complexity.LogProg.compileTM, Complexity.LogProg.ctr, Complexity.LogProg.leftSide, Complexity.LogProg.mkCfg, Complexity.LogProg.mkCfg_state, Complexity.LogProg.nextRet, Complexity.LogProg.retR_scan, Turing.Cfg, Turing.Cfg.workTapeSymbols, Turing.FinTM.bufferTape, Turing.MultiTapeTM.runFrom, Turing.MultiTapeTM.runFrom_succ_eq_step, Turing.MultiTapeTM.step}
\difficulty{6}
Let $P$ be a register program with $m$ registers and $d$ deciders, let $D$ be a Boolean
multi-tape Turing machine with $k_D$ work tapes, and let $x$ be a Boolean input. Suppose
node $l$ has call specification $cs$, and the register selected by argument position $a$
is $rr$. Let register $rr$ contain the finite word $w$ according to the buffer-tape
convention. Assume its initial head position lies between $-1$ and $|w|$, inclusive; if
it is at $-1$, assume the target lies to the left, and if it is at $|w|$, assume the
target lies to the right. Starting in the non-scanning register-return phase for
argument position $a$, there is a time $T$ at which the machine has advanced to the
return successor for argument position $a+1$, with the head of register $rr$ at zero and
all other register head positions unchanged. At every time $t\le T$, each register head
either remains at its initial position or, for register $rr$, lies between $-1$ and the
maximum of its initial position and $0$; the decider tape head positions remain
unchanged throughout.
\end{lemma}

\begin{definition}[Restored register position]
\lean{Complexity.LogProg.retPos}
\label{Complexity.LogProg.retPos}
\leanok
\uses{Complexity.LogProg.CallSpec}
Let a call specification have an ordered list of register indices, and let $rp_0$ assign
an integer position to each register. For a natural number $a$, define the restored
position of register $r$ to be $0$ if $r$ occurs in that list at an index less than $a$;
otherwise define it to be its position $rp_0(r)$.
\end{definition}

\begin{definition}[Register position constraints]
\lean{Complexity.LogProg.RegBox}
\label{Complexity.LogProg.RegBox}
\leanok
\uses{Complexity.LogProg.CallSpec}
Let a call specification list register indices, let $c_0$ assign an integer position to
every register, and let $W$ assign a finite Boolean word to every register. A position
assignment $p$ satisfies the register box if, for every register $r$, it equals $c_0(r)$
when $r$ is absent from the call's list; when $r$ occurs in the list, it lies between
$-1$ and the length of $W(r)$, inclusive.
\end{definition}

\begin{lemma}[Returning after simulating a call]
\lean{Complexity.LogProg.regs_run}
\label{Complexity.LogProg.regs_run}
\leanok
\uses{Complexity.LogProg.CallSpec, Complexity.LogProg.RegBox, Complexity.LogProg.compileTM, Complexity.LogProg.leftSide, Complexity.LogProg.mkCfg, Complexity.LogProg.nextRet, Complexity.LogProg.retPos, Complexity.LogProg.retR_run, Turing.FinTM.bufferTape, Turing.MultiTapeTM.runFrom, Turing.MultiTapeTM.runFrom_add, Turing.MultiTapeTM.runFrom_zero}
\difficulty{6}
Let a register program $P$ call a decider at state $l$, with a call specification whose
argument list has no repetitions. Let $s$ be a segment index, $dir$ and $b$ Boolean
values, and $W$ the words supplied for the registers. Suppose the register tapes store
these words in buffer form for every argument register, their initial head positions lie
in the specified register box, and each argument register at a left endpoint forces the
side predicate to be true while one at a right endpoint forces it to be false. For every
pair of natural numbers $k,a$ with $a+k$ equal to the number of arguments, there is a
time $T$ such that at every time through $T$ the simulated configuration's register-head
positions remain in the register box and its decider-head positions equal the fixed
initial decider-head positions. At time $T$, the compiled machine is in the ordinary
program state corresponding to the call's positive-answer successor if $b$ is true and
its negative-answer successor otherwise, with the original input head, register tapes,
decider tapes, decider-head positions, and output, and with each argument register head
reset according to the return-position rule for the full argument list.
\end{lemma}
